\subsection{System Performance and Scalability}\label{sec:res-live}

\emph{Methodology.} Hour-long live sweeps with up to 32 station processes, a broker, and an edge server on one laptop proved unsuitable for timing: under sustained load the host repeatedly entered power-limited states in which all code ran 4--8$\times$ slower, and the station processes competed for cores that real stations would not share. We therefore measure the service-time distributions of every component in short, speed-checked runs (station model: 300 frames on one core; edge: 60 batches for each $b=1\ldots8$; JPEG payloads of all test frames) and feed them into a trace-driven discrete-event simulator of Fig.~\ref{fig:arch}: periodic capture, a single-core station, a FIFO uplink per station with the emulated rate and delay, a single batching edge server ($b_{\max}=8$), and the downlink delay. About 15\,\% of the edge samples were 5--10$\times$ slower than the median because of the same host throttling; samples above twice the median are excluded, and a slower edge is studied explicitly below. The only free parameter, a constant messaging and serialisation overhead per offloaded frame, is calibrated on one live configuration (Wi-Fi, edge-only, $N=1$): 38.7\,ms.

\emph{Validation.} Table~\ref{tab:validation} compares the simulator with short live runs of the complete system. On Wi-Fi the simulated median latency of the cascade is within 7\,ms of the live value and its p95 within 28\,ms; for edge-only with $N=4$ the live p95 is 60\,ms higher. Deadline-miss rates agree in every Wi-Fi configuration and, on the constrained link, within 15 percentage points. The live edge was 1.4--2.6$\times$ busier than simulated in the unflagged configurations, because on the testbed it shares the CPU with the station processes and the broker; the ``edge 2$\times$ slower'' scenario brackets this case. Four live runs show signs of transient host throttling (marked $\dagger$).

\begin{table}[t]
\centering
\caption{Validation of the trace-driven simulator against the live testbed (live / simulated, both 12\,s at 5 frames/s; simulation: mean of 20 seeds). Latency in ms, misses in \%. $^\dagger$Live run affected by transient host throttling (station median $>2\times$ its trace median, or edge utilisation $>3\times$ simulated).}
\label{tab:validation}
\footnotesize
\setlength{\tabcolsep}{3pt}
\begin{tabular}{@{}llccccc@{}}
\toprule
Link & Policy & $N$ & p50 & p95 & Miss & Edge util. \\
\midrule
wifi$^\dagger$ & station-only & 1 & 143 / 36 & 161 / 40 & 0.0 / 0.0 & 0.00 / 0.00 \\
wifi & station-only & 4 & 53 / 36 & 75 / 40 & 0.0 / 0.0 & 0.00 / 0.00 \\
wifi & cascade & 1 & 115 / 108 & 143 / 115 & 0.0 / 0.0 & 0.16 / 0.06 \\
wifi & cascade & 4 & 94 / 99 & 142 / 115 & 0.0 / 0.0 & 0.24 / 0.18 \\
wifi & edge-only & 1 & 73 / 73 & 80 / 78 & 0.0 / 0.0 & 0.20 / 0.08 \\
wifi & edge-only & 4 & 96 / 73 & 143 / 83 & 0.0 / 0.0 & 0.55 / 0.32 \\
constrained & station-only & 1 & 46 / 36 & 153 / 40 & 0.0 / 0.0 & 0.00 / 0.00 \\
constrained$^\dagger$ & station-only & 4 & 170 / 36 & 191 / 40 & 0.8 / 0.0 & 0.00 / 0.00 \\
constrained$^\dagger$ & cascade & 1 & 678 / 402 & 807 / 498 & 95.0 / 80.5 & 0.61 / 0.06 \\
constrained$^\dagger$ & cascade & 4 & 679 / 310 & 1482 / 447 & 60.0 / 56.1 & 0.96 / 0.18 \\
constrained & edge-only & 1 & 683 / 667 & 1016 / 953 & 100.0 / 100.0 & 0.15 / 0.07 \\
constrained & edge-only & 4 & 1095 / 290 & 3556 / 855 & 100.0 / 100.0 & 0.80 / 0.31 \\
\bottomrule
\end{tabular}
\end{table}


\begin{figure*}[t]
\centering
\includegraphics[width=0.92\textwidth]{fig_sim.png}
\caption{Validated trace-driven simulation on the Wi-Fi link (10 seeds $\times$ 60\,s per point): p95 decision latency (dashed line: 200-ms frame period), deadline misses, and edge utilisation versus the number of stations. Dotted: edge server twice as slow; thin dashed: load model (\ref{eq:load}) with $\mu(8)$.}
\label{fig:sim}
\end{figure*}

\emph{Scaling on Wi-Fi} (Fig.~\ref{fig:sim}). Station-only inspection scales trivially---each decision takes 40\,ms (p95) on the station's own core and no uplink is used---but inherits the station model's accuracy (F1 0.659). Edge-only inspection meets the 200-ms deadline for up to $N=16$ stations (p95 154\,ms, 0.3\,\% misses), at which point the edge is saturated ($\rho=0.99$); at $N=24$ virtually every decision is late. The cascade supports $N=24$ stations (p95 155\,ms, 0.2\,\% misses, $\rho=0.87$) and degrades gracefully: at $N=32$ only 13.6\,\% of decisions miss the deadline, and even in overload at most the offloaded half is late, because clear cases are still decided locally. With an edge twice as slow, the cascade supports 8 stations and edge-only 4. The cascade thus raises the number of stations per edge server by 1.5--2$\times$ while delivering 99\,\% of the edge-only F1, and reduces the mean uplink per station from 821 to 504\,kbit/s (39\,\%). The load model (\ref{eq:load}) with $\mu(8)=96$ frames/s predicts saturation at $N_{\max}=\mu/(fo)\approx19$ stations for edge-only and $\approx37$ for the cascade ($o=0.52$), consistent with the simulated knees between 16--24 and 24--32 stations; at light load it underestimates utilisation because batches are small ($\mu(1)=64$). It is a useful first-order sizing rule: one CPU edge server per $\approx\!\mu/(fo)$ stations.

\emph{Constrained uplinks.} At 1\,Mbit/s a single 14--26\,kB offload needs 110--210\,ms of transmission, which alone exhausts the 200-ms budget. Edge-only decisions are therefore always late (100\,\% misses for every $N$), the cascade misses essentially its offloaded fraction (51--53\,\% for $N\ge8$), and only station-only inspection meets the deadline. On such links the station model must take all time-critical decisions; offloaded verdicts remain useful as delayed quality records, or the payload must shrink (smaller crop, stronger compression).
